\begin{frame}{(d) 추천: 사용자-아이템도 그래프다}
  \begin{itemize}
    \item 넷플릭스/아마존이 푸는 문제: ``이 사용자가 다음에 뭘 보고
      싶어 할까?''
    \item \textbf{사용자 노드}와 \textbf{아이템 노드}를 각각 두고,
      ``봤음/샀음/평점을 매김''을 \textbf{엣지}로 놓으면 ---
      \kb{이분 그래프}(bipartite graph): 사용자끼리, 아이템끼리는
      직접 연결되지 않는다.
    \item (a)/(b)/(c)와 똑같은 패턴: \textbf{그래프를 어떻게
      만드는가}(무엇을 노드로, 무엇을 엣지로)가 추천 품질을
      결정한다.
  \end{itemize}
  \vskip0.6em
  \begin{block}{13.3절 개인화 PageRank가 이미 이 문제를 가리켰다}
    ``무작위 걷기를 주어진 노드에서 다시 시작하는 변형 --- 추천
    시스템에 사용'' --- 여기서 그 문장을 직접 풀어본다.
  \end{block}
\end{frame}
